\section{Implementation}
\subsection{Platform and Tool Flow}
The prototype is a 2$\times$3 ESP SoC on a VCU118 FPGA, with an Ariane RISC-V processor, NVDLA \texttt{nv\_small}, and a SAT-Accel tile. The nominal tile clock is 78.125\,MHz; per-run counter calibration is approximately 78.1 cycles per microsecond. Linux hosts the application and accelerator drivers. Both accelerator paths use coherent DMA.

The original project proposal named Chipyard and Gemmini. The implemented system uses ESP and NVDLA to match the available integration and tools. SAT-Accel's seven HLS kernels are built with Vitis HLS 2023.2. ReasonLink's hardware delta generator is implemented directly in Verilog RTL; the CPU-generated delta path is software executed by Ariane. On-chip base reuse is implemented in the HLS solver and is available to both delta-generation paths. A custom block design, APB-to-AXI-Lite shell, Linux driver, and user library replace the original \texttt{v++}/XRT integration. Floating-point VSIDS operations remain in the solver, and clause capacity is reduced for the target FPGA.

The resident-base out-of-context SAT-Accel netlist uses 113,062 LUTs; adding ReasonLink increases this to 117,698 LUTs, an incremental 4,636 LUTs (approximately 4.1\%). Both configurations use 434 URAMs, approximately 45.2\% of the device. These are out-of-context synthesized resources, not isolated placed-and-routed overhead. The November 4 rebuilt full-system bitstreams meet timing, with recorded worst setup slack of $+$0.774\,ns for the SAT-Accel build and $+$0.369\,ns for the ReasonLink build.

\draftnote{Archive the final utilization and timing reports alongside exact bitstream and source hashes. Distinguish final routed resources from HLS estimates and out-of-context synthesis in the final resource table. Support any energy-efficiency claims with energy measurements.}

\draftnote{Rebuild the Linux image with the output-layout fix and package the exact executables, bitstreams, source revisions, raw results, and figure-generation scripts used for the final evaluation. Condense the integration bug history to the details needed to explain the design and its validation.}

\subsection{Classifier and Workload}
We retain ESP's LeNet architecture and retrain weights with preprocessing compatible with the int8 runtime. The recorded MNIST test accuracies are 99.17\% in floating point and 99.20\% with the int8 emulator. The board loadable exposes the logits used by the handoff; softmax is unnecessary for argmax.

The generated dataset contains 1,000 unique-solution puzzles with 24--40 clues, with a mean of approximately 32.1. Each given is paired with an MNIST test image. The current board study uses 100 puzzles. Misclassification can make a previously unique puzzle unsatisfiable or introduce alternative solutions, so solver correctness must be assessed against predicted givens, separately from recovery of the original grid.

\subsection{Integration Validation}
Integration exposed failures that component-level tests did not capture. ESP's APB reads could arrive without the setup phase expected by our bridge. ESP's AXI socket also ignored write strobes, allowing result writes to overwrite neighboring bytes; a wrapper adapter corrected the behavior. Full-SoC simulation and a bare-metal board solve validated the resulting path.

On Linux, interrupt acknowledgement and overlapping accelerator DMA reservations caused additional failures. These were corrected before collecting the accepted comparisons. Finally, ReasonLink initially decoded digits 8 and 9 from incorrect channels because the runtime's output stride could not be represented by the selected layout encoding. The driver now uses the supported contiguous layout and rejects unrepresentable strides. Incorrect pre-fix CSVs are preserved for debugging but excluded from the reported result set.
